The grouped body is decoded and each body instruction is checked for REPG repeatability before execution begins.
REPG has no entry zero-count annul: after a valid group has passed those checks, its first iteration begins even
when the selected loop-control register is zero.

The selected Rn is the live unsigned loop-control value; REPG does not capture a separate remaining count. Body
instructions read and write that register by their ordinary architectural rules. After one whole grouped-body
iteration succeeds, REPG examines the value then held in the selected register. If it is zero, REPG leaves it
unchanged and exits repeat state. Otherwise, REPG first decrements it by one and then uses the resulting value to
select the next action: zero exits repeat state, while nonzero continues at the group start. This guarded decrement
and resulting-value decision form the group-boundary state transition. With no body write to the register, an
initial zero executes one whole-group iteration, while an initial positive value N executes exactly N iterations;
both cases leave zero. A body write affects the boundary reached by that same iteration: writing zero exits without
a decrement, and writing one decrements to zero and exits. The REPGF assembler-checked alias rejects bodies that
write the selected loop-control register.

Completed prior body instructions in a faulting iteration remain committed, and later body instructions do not
execute. The group-boundary state transition does not occur for an incomplete iteration. The saved repeat
continuation records the selected loop-control register, condition T, repeat kind,
group\_start\_delta, and body\_bytes in FRAME\_EXT1. FLAGS and FFLAGS follow the rules of the last completed grouped
instruction; FFLAGS from completed floating-point operations are accumulated by bitwise inclusive-or.
